%% Chapter 9 - Conclusion and Future Directions
%% Sections: 9.1 Conclusion; 9.2 Achievement of objectives; 9.3 Individual
%% research and product contributions; 9.4 Future directions.
%% Objective mapping uses only achieved/partially achieved/not-achieved
%% statuses backed by evidence.

\chapter{Conclusion and Future Directions}\label{ch:conclusion}

This chapter closes the dissertation. It restates the central conclusion,
accounts for each objective against the experiment record with an explicit
status, isolates the individual research and product contributions, and
itemises the directions that follow from the limitations committed in
Chapter~\ref{ch:discussion}.

\section{Conclusion}\label{sec:conclusion}

The central question was whether a small, offline, unattended language model
can be made reliable enough to drive an automated security copilot once its
output is subject to deterministic validation. The measured answer is yes,
within a closed threat model, on the evidence of this dissertation.

On a committed 1,500-record population, the guarded pipeline raised the
flaw-catch rate from 0\% to 64--72\%, lifted decision accuracy from 69--78\%
to 82--86\%, and halved the false-positive rate to 14--15\% on two distinct
offline backends (Section~\ref{sec:exp1}). The ablation record localised the
cause: the deterministic validation stack is the dominant safety contributor
($-30$ and $-24$ catch points when removed), the hallucination validators are
the largest single-layer contributors, and the per-class catch spectrum tracks
the decidability of the underlying predicate (Sections~\ref{sec:exp3} and
\ref{sec:exp4}). Safety generalises across backends while individual accuracy
contributions do not (Section~\ref{sec:exp2}), which makes deterministic
guardrailing\emph{---}rather than model choice\emph{---}the primary driver of
the measured safety.

The conclusion is therefore not that the engine eliminates model error. It is
that the engine converts an open problem---making a language model
trustworthy---into the closed problem of classifying a single utterance under a
known policy, and that the deterministic half of that reduction does most of
the safety work. Three boundaries bound the claim and are restated because they
are the honest frame for the conclusion: the evaluation is single-run point
estimates without confidence intervals, the data set is synthetic and
non-uniform, and the threat model closes at retrieval-free, no-free-form-output
boundary (Sections~\ref{sec:discussion-limitations} and \ref{sec:ethics}).

\section{Achievement of Objectives}\label{sec:achievement-objectives}

The four objectives of Section~\ref{sec:aim-and-objectives} are accounted for
individually below. Each status is grounded in the cited evidence; no objective
is marked achieved on assertion alone.

\begin{itemize}
\item \textbf{Objective 1\emph{---}design and implement the Intent Recognition
Engine}\textnormal{: a component that receives free-text security requests and
emits strictly schema-conformant, authorised intent contracts over a small
offline language model.}
\textbf{Status: achieved.} The seven-stage core and five middleware components
are implemented under \texttt{src/} and exercised by 506 automated tests (76\%
statement coverage; 434 unit and 72 integration) described in
Section~\ref{sec:component-implementation} and
Section~\ref{sec:product-results}. The live C1--C2 contract verification
confirmed that real code paths emit and consume validated,
schema-conformant contracts (Section~\ref{sec:integration-results}).

\item \textbf{Objective 2\emph{---}construct a purpose-built evaluation data set
of security prompts, covering well-formed requests and deliberately flawed
prompts across the flaw categories relevant in practice, each annotated with
gold labels.}
\textbf{Status: achieved.} The 1,500-record set (600 well-formed, 900 flawed)
spans the eight+ flaw classes used in the experiments, is versioned with the
project source, and underpins every measurement in
Chapter~\ref{ch:results} (Section~\ref{sec:data-materials}).

\item \textbf{Objective 3\emph{---}establish an evaluation protocol that
quantifies the pipeline improvement over raw output and abates each guard layer
independently.}
\textbf{Status: achieved.} The protocol is committed in advance, applied to the
fixed data set, and recorded in a manifest for independent audit
(Section~\ref{sec:reproducibility}). Experiments 1 and 3 operationalise the
raw-versus-guarded comparison and the per-layer ablation respectively
(Sections~\ref{sec:exp1} and \ref{sec:exp3}).

\item \textbf{Objective 4\emph{---}compare two offline language-model backends
under the identical guardrail stack.}
\textbf{Status: achieved.} Both backends were evaluated under the identical
stack; the safety ordering was consistent across backends while per-layer
accuracy contributions differed, a split recorded in Sections~\ref{sec:exp2}
and \ref{sec:exp3} and analysed against RQ3 in
Section~\ref{sec:discussion-interpretation}.
\end{itemize}

All four objectives are thus achieved. Because the evaluation is single-run and
the three semantic flaw classes remain weakly protected, achievement is on the
evidence available within the scope committed in Section~\ref{sec:scope-limitations-assumptions};
the boundary conditions that would warrant a ``partially achieved''
qualification are the same boundaries the limitations section records
(Section~\ref{sec:discussion-limitations}).

\section{Individual Research and Product Contributions}\label{sec:contributions}

\subsection*{Research contributions}

The research contribution is a measured result with three parts. First, the
deterministic validation stack is shown to be the dominant driver of the
measured safety and is reproducible across two backends: no single layer, and
no model, produces the catch rate the composed stack produces
(Section~\ref{sec:exp3}). Second, the catch effectiveness boundary is
predictable from decidability: per-class catch runs from 100\% on structurally
decidable classes down to 40\% on meaning-dependent classes, which converts the
limitation of single-shot offline inference into a design law---ask a model for
structure, never for a verdict (Section~\ref{sec:exp4}). Third, safety and
usability improve together: the guarded pipeline raises accuracy, raises catch
and halves false positives simultaneously (Section~\ref{sec:exp1}), weakening
the implicit assumption that deterministic checks must trade precision for
recall or fluency for utility.

\subsection*{Product contributions}

The product contribution is the IRE itself: a shippable,
deployment-gated, offline component whose output is machine-consumed, not
reviewed---the empty corner of the existing-systems survey (Section~\ref{sec:existing-systems}
and Table~\ref{tab:existing-systems}). The product properties most worth
exporting elsewhere are architectural: the safety invariant
(Equation~\ref{eq:safety-invariant}), the audit-friendly structured response,
and the executable test suite of 506 tests that encodes the invariant
(Section~\ref{sec:product-results}). The verified integration findings
(Section~\ref{sec:integration-results}) confirm that when the engine is asked
to produce a contract for a downstream component, the contract it produces is
validated, in-scope and authorised on the code paths that would run in
production.

Both contributions remain subject to the honest engineering caveats itemised in
Section~\ref{sec:discussion-product}: several documented dependencies are
doc-only, the default development API key bypasses transport authentication if
left untouched, and the \texttt{production} scope-mode value is dead
configuration that falls through to \texttt{WARN}. These are stated as
product-integrity gaps to close, and they do not affect the validity of the
evaluation, which ran with production credentials and strict scope semantics.

\section{Future Directions}\label{sec:future-directions}

The future directions follow directly from the limitations and open items
committed in Chapters~\ref{ch:results} and \ref{ch:discussion}. Each is listed
with the state of the evidence that motivates it.

\emph{Remove the single-run limitation.} Every configuration in
Experiments 1--4 was executed once (Section~\ref{sec:procedure}). Repeating the
committed protocol with seed sweeps and confidence intervals would convert the
point estimates into bounded estimates and would resolve whether the
negative-$\Delta$Catch ablation entries are stable (Section~\ref{sec:exp3}).

\emph{Close the semantic-class gap.} The three meaning-dependent flaw classes
(contradictory, out-of-scope-target, ambiguous) remain caught at 40--59\%
because no deterministic signal exists for them (Section~\ref{sec:exp4}).
Structured prompting that forces the model to enumerate decision grounds,
targeted enumerators for scope, or a bounded-sampling escalator on the
agreement gate are concrete candidate directions within the offline budget
(Section~\ref{sec:related-research}).

\emph{Complete the production validation.} The C1--C2 happy-path re-run against
the production \texttt{gemma4:e4b} configuration is recorded as a TODO in
Section~\ref{sec:integration-results} until the model is provisioned on the
evaluation machine. Latency, throughput and energy profiling on the target
hardware were not measured and are required before interactive-use claims
(Section~\ref{sec:discussion-limitations}).

\emph{Close the engineering configuration gaps.} Shipping hardening requires
removing the default development API key, wiring the \texttt{production}
scope-mode value to the \texttt{BLOCK} path of the enforcement policy, and
reconciling the doc-only dependencies (spaCy, jsonschema runtime, LoRA,
Mahalanobis and XAI feature-flag paths) with their documented roles
(Section~\ref{sec:discussion-product}).

\emph{Extend the evaluation surface.} The threat model closes at the
retrieval-free, no-free-form-output boundary (Section~\ref{sec:ethics}). Future
work could expand the boundary to indirect prompt injection, device-compromise
and backdoored-weights scenarios, or relax the deployment constraint to measure
what a fine-tuned backend (the planned LoRA Phase~2 work) adds to the safety
record.

\emph{Re-verify the commercial position.} The comparative assessment cites no
named commercial product for lack of verifiable primary sources
(Section~\ref{sec:existing-systems}); a commercialisation step should survey
named products with primary sources and re-verify the gap claim before any
positioning statement is made (Section~\ref{sec:market-position}).